\chapter{Weak Convergence and Donsker's Theorem}
\begin{center}\itshape
Take any random walk whose steps have mean $0$ and variance $1$, and look at its first $n$ steps from far away,\\
shrinking time by $n$ and space by $\sqrt n$. As $n$ grows, the picture becomes Brownian motion,\\
whatever the law of the steps. So any question about the walk at large scale has the answer it has for Brownian motion.
\end{center}

In 1946 Paul Erd\H{o}s and Mark Kac wanted the limiting law of the maximum of the partial sums $S_1,\dots,S_n$ of i.i.d.\ variables. They
noticed something odd: the limit did not depend on the law of the steps, as long as the variance was finite. So they could compute it for
the simplest walk, the coin-tossing walk, where combinatorics gives exact formulas, and then transfer the answer to every other walk. They
called this an \emph{invariance principle}, and they did the same for the number of positive partial sums, recovering L\'evy's arcsine
law. Their proofs were clever but special: each statistic needed its own argument.

Monroe Donsker saw what was behind the trick. In a memoir of 1951, drawn from his thesis, he proved that the whole rescaled path of the walk, as a random continuous
function, converges in distribution to Brownian motion. Then every statistic that is a continuous function of the path converges, at
once, to the same statistic of Brownian motion: the maximum, the time spent positive, the time of the last zero, and any other you care to
name. Two years earlier Joseph Doob had made the same observation heuristically for the Kolmogorov--Smirnov statistic of goodness of fit,
whose limit law Kolmogorov had computed in 1933 by a long calculation; Donsker's theorem for empirical processes, in 1952, made Doob's
heuristic a proof. Since then ``Donsker'' has been the name of a whole industry: the functional limit theorems of probability and
statistics.

This chapter proves Donsker's theorem in the way discovered by Anatoliy Skorokhod in 1961, which avoids the general theory of tightness in
spaces of functions. Skorokhod showed that any centred law with a finite variance can be realised as the position of a Brownian motion at a
well-chosen stopping time. Running Brownian motion and stopping it again and again produces a random walk with any step law you like, sitting
inside the Brownian path. Then the walk and the Brownian motion are close, path by path, and Donsker's theorem follows from the strong law
of large numbers. The same embedding gives the \emph{law of the iterated logarithm} of Philip Hartman and Aurel Wintner (1941): the walk
reaches $\sqrt{2n\log\log n}$ infinitely often and exceeds it by any fixed factor only finitely often, a statement that Khintchine had found in
1924 for coin tossing.

Before that we need the language: weak convergence in $\R^d$, where L\'evy's continuity theorem turns characteristic functions into limit
theorems, and weak convergence in the space $C[0,1]$ of continuous paths.

\begin{reminder}[from Tomes 3 and 6]
\emph{Weak convergence} (Tome~6, Chapter~11): probabilities $\mu_n$ on a metric space converge weakly, $\mu_n\Rightarrow\mu$, if $\int f\dd\mu_n\to\int f\dd\mu$
for bounded continuous $f$. By the \emph{portmanteau theorem} (Tome~6, Theorem~11.4) it is enough to test bounded Lipschitz $f$, or to check
$\limsup\mu_n(F)\le\mu(F)$ for closed $F$. A sequence of probabilities on $\R^d$ is \emph{tight} if for every $\varepsilon$ some compact set has mass
$\ge1-\varepsilon$ under all of them; \emph{Prokhorov's theorem} (Tome~6, Theorem~11.7): a tight sequence on $\R^d$ has a weakly convergent subsequence.
\emph{Brownian motion} (Tome~3, Chapter~11): $B_0=0$, continuous paths, independent increments $B_t-B_s\sim N(0,t-s)$; it exists (Theorem~11.4); its
finite-dimensional laws are centred Gaussian with covariance $\min(s,t)$ (Proposition~11.2); $c^{-1}B_{c^2t}$ and $B_{s+t}-B_s$ are Brownian motions, the
latter independent of $\mathcal F_s=\sigma(B_r,r\le s)$ (Proposition~11.5); $B_t$ and $B_t^2-t$ are martingales (Proposition~11.11); the maximum
$M_t=\max_{s\le t}B_s$ has the law of $|B_t|$ (Theorem~11.12). \emph{Simple random walk} (Tome~3, Chapter~7): $u_{2n}=P(S_{2n}=0)$, the first-return
decomposition $u_{2n}=\sum_{r\le n}f_{2r}u_{2n-2r}$ (Lemma~7.2), $P(S_1\ge0,\dots,S_{2n}\ge0)=u_{2n}$ (Theorem~7.3), and the arcsine limit of
$\sum_{an\le k\le bn}u_{2k}u_{2n-2k}$ (Theorem~7.7).
\end{reminder}

\section{L\'evy's continuity theorem}

For a random vector $X$ in $\R^d$, the characteristic function is $\varphi_X(t)=\mathbb E e^{i\langle t,X\rangle}$, $t\in\R^d$. Tome~3 proved that it
determines the law in dimension~$1$. In any dimension, a Gaussian smoothing does it.

\begin{lemma}[uniqueness]\label{lem:chfunique}
Two random vectors of $\R^d$ with the same characteristic function have the same law.
\end{lemma}
\begin{proof}
Let $Z$ be a standard Gaussian vector of $\R^d$ independent of $X$, and $\sigma>0$. The density of $\sigma Z$ is
$g_\sigma(y)=(2\pi)^{-d}\int e^{-i\langle t,y\rangle}e^{-\sigma^2|t|^2/2}\dd t$ (Fourier inversion of the Gaussian characteristic function, coordinate by
coordinate, Tome~3, Theorem~6.8). So $X+\sigma Z$ has the density
\[
p_\sigma(x)=\mathbb E\,g_\sigma(x-X)=(2\pi)^{-d}\int e^{-i\langle t,x\rangle}\varphi_X(t)\,e^{-\sigma^2|t|^2/2}\dd t,
\]
by Fubini, which applies since the integrand is bounded by $e^{-\sigma^2|t|^2/2}$. Thus $\varphi_X$ determines the law of $X+\sigma Z$ for every $\sigma>0$.
As $\sigma\to0$, $X+\sigma Z\to X$ surely, hence in distribution, and weak limits are unique. So $\varphi_X$ determines the law of $X$.
\end{proof}

The next inequality turns the behaviour of $\varphi$ near $0$ into control of the tails.

\begin{lemma}[truncation inequality]\label{lem:truncineq}
For a real random variable $X$ and $u>0$,
\[
P\Bigl(|X|\ge\frac2u\Bigr)\le\frac1u\int_{-u}^u\bigl(1-\varphi_X(t)\bigr)\dd t .
\]
\end{lemma}
\begin{proof}
For real $x\ne0$, $\frac1{2u}\int_{-u}^u(1-e^{itx})\dd t=1-\frac{\sin ux}{ux}$, which is $\ge0$ always and $\ge\frac12$ when $|ux|\ge2$, since then
$|\frac{\sin ux}{ux}|\le\frac12$; for $x=0$ it is $0$. Take expectations and use Fubini: $\frac1{2u}\int_{-u}^u(1-\varphi_X(t))\dd t\ge\frac12P(|uX|\ge2)$.
\end{proof}

\begin{theorem}[L\'evy's continuity theorem]\label{thm:levycont}
Let $X_n$ be random vectors of $\R^d$ with characteristic functions $\varphi_n$. If $\varphi_n(t)\to\varphi(t)$ for every $t$, and $\varphi$ is continuous at
$0$, then $\varphi$ is the characteristic function of a random vector $X$ and $X_n\Rightarrow X$. Conversely, $X_n\Rightarrow X$ implies $\varphi_n\to\varphi_X$
pointwise.
\end{theorem}
\begin{proof}
The converse holds because $x\mapsto e^{i\langle t,x\rangle}$ is bounded and continuous.

\emph{Tightness.} Let $\varepsilon>0$ and let $e_j$ be the $j$-th basis vector. The coordinate $X_n^{(j)}$ has characteristic function $s\mapsto\varphi_n(se_j)$.
Since $\varphi(0)=\lim\varphi_n(0)=1$ and $\varphi$ is continuous at $0$, there is $u>0$ with $\frac1u\int_{-u}^u|1-\varphi(se_j)|\dd s<\varepsilon$ for all $j$. By
dominated convergence ($|1-\varphi_n|\le2$), $\frac1u\int_{-u}^u(1-\varphi_n(se_j))\dd s\to\frac1u\int_{-u}^u(1-\varphi(se_j))\dd s$, so by
Lemma~\ref{lem:truncineq}, $P(|X_n^{(j)}|\ge2/u)<2\varepsilon$ for all $n\ge n_0$ and all $j$. Enlarging the box $[-2/u,2/u]^d$ to cover the finitely many
$n<n_0$ as well, we get a compact set of mass at least $1-2d\varepsilon$ under every law of $X_n$.

\emph{Identification.} By Prokhorov, every subsequence of $(X_n)$ has a further subsequence converging in distribution to some $Y$; then $\varphi_Y=\lim\varphi_n
=\varphi$ along it. So $\varphi$ is a characteristic function, of $X$ say, and by Lemma~\ref{lem:chfunique} every such limit $Y$ has the law of $X$. If
$X_n\not\Rightarrow X$, there would be a bounded continuous $f$, an $\eta>0$ and a subsequence with $|\mathbb E f(X_n)-\mathbb E f(X)|\ge\eta$; a further
subsequence converges to the law of $X$, a contradiction.
\end{proof}

\begin{corollary}[Cram\'er--Wold device]\label{cor:cramerwold}
$X_n\Rightarrow X$ in $\R^d$ if and only if $\langle t,X_n\rangle\Rightarrow\langle t,X\rangle$ for every $t\in\R^d$.
\end{corollary}
\begin{proof}
$\varphi_{X_n}(t)=\varphi_{\langle t,X_n\rangle}(1)$. If the projections converge, the characteristic functions converge pointwise to $\varphi_X$, which is
continuous; the theorem applies. The converse is the continuous mapping $x\mapsto\langle t,x\rangle$.
\end{proof}

So a multivariate limit theorem reduces to one-dimensional ones, and in particular the central limit theorem for i.i.d.\ vectors with a covariance
$\Sigma$ follows from Tome~3's: $\frac1{\sqrt n}\sum_{k\le n}(X_k-\mathbb E X_k)\Rightarrow N(0,\Sigma)$, because each projection $\langle t,\cdot\rangle$ is a real
CLT with variance $t^{\mathsf T}\Sigma t$, and $N(0,\Sigma)$ has the characteristic function $e^{-t^{\mathsf T}\Sigma t/2}$ (Tome~3, Theorem~6.15). Chapter~4
returns to this with rates.

\section{Random continuous functions}

Let $C=C[0,1]$ with the norm $\|f\|=\max_{t\le1}|f(t)|$. It is a complete separable metric space: polynomials with rational coefficients are dense by
Weierstrass's theorem.

\begin{lemma}\label{lem:cborel}
The Borel $\sigma$-algebra of $C$ is generated by the evaluations $\pi_t:f\mapsto f(t)$, $t\in[0,1]$. Consequently a process $(X_t)_{t\le1}$ with continuous paths
is a random element of $C$, and its law on $C$ is determined by its finite-dimensional distributions.
\end{lemma}
\begin{proof}
Each $\pi_t$ is continuous, hence Borel. Conversely, the closed ball $\{f:\|f-g\|\le r\}$ equals $\bigcap_{t\in\Q\cap[0,1]}\{f:|f(t)-g(t)|\le r\}$ by continuity,
so it lies in $\sigma(\pi_t)$. Every open set of the separable space $C$ is a countable union of such balls (with centres in a countable dense set and rational
radii), so open sets, hence all Borel sets, lie in $\sigma(\pi_t)$. If $X$ has continuous paths, $\{X\in A\}$ is an event for $A$ in the generating family
$\{\pi_t^{-1}(B)\}$, hence for all Borel $A$. Finally, the sets $\{f:(f(t_1),\dots,f(t_k))\in B\}$ form a $\pi$-system generating the Borel sets, and two laws
agreeing on it agree everywhere (Tome~6, Theorem~6.5).
\end{proof}

Two tools transfer weak convergence. For $h:S\to S'$ between metric spaces, let $D_h$ be its set of discontinuity points.

\begin{proposition}[continuous mapping]\label{prop:contmap}
If $X_n\Rightarrow X$ in $S$, $h:S\to S'$ is measurable, and $D_h$ is contained in a Borel set $N$ with $P(X\in N)=0$, then $h(X_n)\Rightarrow h(X)$.
\end{proposition}
\begin{proof}
Let $F\subseteq S'$ be closed. If $x$ is in the closure of $h^{-1}(F)$ and $h$ is continuous at $x$, then $h(x)\in F$; so $\overline{h^{-1}(F)}\subseteq h^{-1}(F)\cup N$. By
the portmanteau theorem,
\[
\limsup_nP(h(X_n)\in F)\le\limsup_nP\bigl(X_n\in\overline{h^{-1}(F)}\bigr)\le P\bigl(X\in\overline{h^{-1}(F)}\bigr)\le P(h(X)\in F).
\]
\end{proof}

\begin{lemma}[converging together]\label{lem:together}
Let $X_n,Y_n$ be random elements of a metric space $(S,d)$ on the same probability space, with $Y_n\Rightarrow Y$ and $d(X_n,Y_n)\to0$ in probability. Then
$X_n\Rightarrow Y$.
\end{lemma}
\begin{proof}
For $f$ bounded and $L$-Lipschitz and $\varepsilon>0$, $|\mathbb E f(X_n)-\mathbb E f(Y_n)|\le L\varepsilon+2\|f\|_\infty P(d(X_n,Y_n)>\varepsilon)$, whose $\limsup$ is
$L\varepsilon$. So $\mathbb E f(X_n)\to\mathbb E f(Y)$ for every bounded Lipschitz $f$, which suffices by the portmanteau theorem.
\end{proof}

\section{The strong Markov property and Skorokhod's embedding}

Let $B$ be a Brownian motion, and let $\mathcal G_t=\sigma(\mathcal G_0,B_s:s\le t)$, where $\mathcal G_0$ is a $\sigma$-algebra independent of $B$ (it will carry
auxiliary randomness). Then $B_{s+t}-B_s$ is independent of $\mathcal G_s$: it is independent of $\sigma(B_r,r\le s)$ by Tome~3, Proposition~11.5, and the pair is
independent of $\mathcal G_0$ (grouping, Proposition~\ref{prop:grouping}). A \emph{stopping time} is a random $\tau\in[0,\infty]$ with $\{\tau\le t\}\in\mathcal G_t$ for
all $t$, and $\mathcal G_\tau$ is the family of events $A$ with $A\cap\{\tau\le t\}\in\mathcal G_t$ for all $t$.

\begin{theorem}[strong Markov property]\label{thm:strongmarkov}
Let $\tau$ be an almost surely finite stopping time. Then $W_t=B_{\tau+t}-B_\tau$ is a Brownian motion, independent of $\mathcal G_\tau$.
\end{theorem}
\begin{proof}
Let $\tau_m=(\lfloor2^m\tau\rfloor+1)2^{-m}$, the first dyadic time of level $m$ strictly after $\tau$; then $\tau_m\downarrow\tau$. Fix $A\in\mathcal G_\tau$, times
$0\le t_1<\dots<t_k$ and a bounded continuous $g:\R^k\to\R$. For $j\ge1$, the event $\{\tau_m=j2^{-m}\}=\{(j-1)2^{-m}\le\tau<j2^{-m}\}$ lies in
$\mathcal G_{j2^{-m}}$, and so does $A\cap\{\tau<s\}=\bigcup_{q<s,\,q\in\Q}(A\cap\{\tau\le q\})$ for $s=j2^{-m}$. With $W^m_t=B_{\tau_m+t}-B_{\tau_m}$, the Markov
property at the fixed time $s=j2^{-m}$ gives
\[
\begin{gathered}
\mathbb E\bigl[g(W^m_{t_1},\dots,W^m_{t_k})\mathbf 1_A\bigr]=\sum_{j\ge1}\mathbb E\bigl[g(B_{s+t_1}-B_s,\dots)\mathbf 1_{A\cap\{\tau_m=s\}}\bigr]\qq
=\sum_{j\ge1}P(A\cap\{\tau_m=s\})\,\mathbb E\,g(B_{t_1},\dots,B_{t_k})=P(A)\,\mathbb E\,g(B_{t_1},\dots,B_{t_k}).
\end{gathered}
\]
As $m\to\infty$, $W^m_t\to W_t$ by continuity of paths, and bounded convergence gives $\mathbb E[g(W_{t_1},\dots,W_{t_k})\mathbf 1_A]=P(A)\,\mathbb E\,g(B_{t_1},\dots,B_{t_k})$.
With $A=\Omega$, $W$ has the finite-dimensional laws of Brownian motion; it has continuous paths and $W_0=0$, so it is a Brownian motion. And since bounded
continuous $g$ determine laws, $(W_{t_1},\dots,W_{t_k})$ is independent of $\mathcal G_\tau$ for all times, hence $W$ is (Lemma~\ref{lem:pisuffice}).
\end{proof}

We also need optional stopping in continuous time.

\begin{lemma}\label{lem:ostbm}
For a stopping time $\tau$ and $K>0$: $\mathbb E B_{\tau\wedge K}=0$ and $\mathbb E B_{\tau\wedge K}^2=\mathbb E[\tau\wedge K]$.
\end{lemma}
\begin{proof}
Let $\sigma=\tau\wedge K$ and $\sigma_m=(\lfloor2^m\sigma\rfloor+1)2^{-m}$, a stopping time for the discrete filtration $(\mathcal G_{j2^{-m}})_j$, bounded by $K+1$.
The sequences $(B_{j2^{-m}})_j$ and $(B_{j2^{-m}}^2-j2^{-m})_j$ are martingales for this filtration, so by discrete optional stopping (Tome~3, Theorem~10.12),
$\mathbb E B_{\sigma_m}=0$ and $\mathbb E B_{\sigma_m}^2=\mathbb E\sigma_m$. As $m\to\infty$, $B_{\sigma_m}\to B_\sigma$ and $\sigma_m\to\sigma$, with $|B_{\sigma_m}|\le M^*=\max_{s\le K+1}|B_s|$.
Since $M^*\le\max_sB_s+\max_s(-B_s)$ and each maximum has the law of $|B_{K+1}|$, $\mathbb E(M^*)^2\le4(K+1)<\infty$, and dominated convergence concludes.
\end{proof}

\begin{proposition}[exit from an interval]\label{prop:exit}
Let $a<0<b$ and $\tau=\inf\{t:B_t\notin(a,b)\}$. Then $\tau<\infty$ almost surely, $P(B_\tau=b)=\frac{-a}{b-a}$, and $\mathbb E\tau=-ab$.
\end{proposition}
\begin{proof}
$\tau$ is a stopping time: $\{\tau\le t\}=\{\min_{s\le t}B_s\le a\}\cup\{\max_{s\le t}B_s\ge b\}$, and by continuity these extrema are attained over rational $s$ and $s=t$. It is
finite because the independent increments $B_{k+1}-B_k$ exceed $b-a$ infinitely often. By Lemma~\ref{lem:ostbm}, $\mathbb E B_{\tau\wedge K}=0$ and $\mathbb E[\tau\wedge K]=\mathbb E
B_{\tau\wedge K}^2$. Let $K\to\infty$: $|B_{\tau\wedge K}|\le\max(-a,b)$, so $\mathbb E B_\tau=0$, which gives $bP(B_\tau=b)+a(1-P(B_\tau=b))=0$; and by monotone convergence on the left,
$\mathbb E\tau=\mathbb E B_\tau^2=b^2\frac{-a}{b-a}+a^2\frac b{b-a}=-ab$.
\end{proof}

\begin{theorem}[Skorokhod's embedding, 1961]\label{thm:skorokhod}
Let $X$ be a random variable with $\mathbb E X=0$ and $\mathbb E X^2<\infty$. There is a stopping time $\tau$ (for a filtration $\mathcal G_t$ as above, with suitable
$\mathcal G_0$) such that $B_\tau$ has the law of $X$ and $\mathbb E\tau=\mathbb E X^2$.
\end{theorem}
\begin{proof}
Let $\mu$ be the law of $X$, and $c=\mathbb E X^+=\mathbb E X^-$; if $c=0$ then $X=0$ and $\tau=0$ works. Otherwise let $(U,V)$ be a random pair independent of $B$, generating
$\mathcal G_0$, with law
\[
P\bigl((U,V)\in\dd u\,\dd v\bigr)=\frac{(v-u)}c\,\mu(\dd u)\,\mu(\dd v)\quad(u<0<v),\qquad P\bigl((U,V)=(0,0)\bigr)=\mu(\{0\}).
\]
The total mass is
\[
\begin{gathered}
\frac1c\Bigl(\mu(0,\infty)\int_{u<0}(-u)\,\mu(\dd u)+\mu(-\infty,0)\int_{v>0}v\,\mu(\dd v)\Bigr)+\mu(\{0\})\\
=\mu(\R\setminus\{0\})+\mu(\{0\})=1.
\end{gathered}
\]
Let $\tau$
be the exit time of $(U,V)$, with $\tau=0$ when $U=V=0$. It is a stopping time for $(\mathcal G_t)$. Given $(U,V)=(u,v)$, Proposition~\ref{prop:exit} (applied to $B$, which is
independent of $\mathcal G_0$) says $B_\tau=v$ with probability $\frac{-u}{v-u}$. So for a Borel set $A\subseteq(0,\infty)$,
\[
P(B_\tau\in A)=\int_{v\in A}\int_{u<0}\frac{v-u}c\cdot\frac{-u}{v-u}\,\mu(\dd u)\,\mu(\dd v)=\mu(A)\cdot\frac1c\int_{u<0}(-u)\,\mu(\dd u)=\mu(A),
\]
and symmetrically for $A\subseteq(-\infty,0)$; and $P(B_\tau=0)=\mu(\{0\})$. Finally
\[
\begin{gathered}
\mathbb E\tau=\mathbb E[-UV]=\frac1c\int_{v>0}\int_{u<0}(-uv)(v-u)\,\mu(\dd u)\,\mu(\dd v)\qq=\int_{v>0}v^2\mu(\dd v)+\int_{u<0}u^2\mu(\dd u)=\mathbb E X^2.
\end{gathered}
\]
\end{proof}

\begin{corollary}[embedding a random walk]\label{cor:embedwalk}
Let $X_1,X_2,\dots$ be i.i.d.\ with mean $0$ and variance $\sigma^2$. On a suitable probability space there are a Brownian motion $B$ and stopping times
$0=T_0\le T_1\le T_2\le\cdots$ such that $(B_{T_n})_{n\ge0}$ has the law of $(S_n)_{n\ge0}$, and the increments $T_n-T_{n-1}$ are i.i.d.\ with mean $\sigma^2$.
\end{corollary}
\begin{proof}
Take a Brownian motion $B$ and an independent i.i.d.\ sequence of pairs $(U_k,V_k)$ with the law of the proof above, and $\mathcal G_0=\sigma((U_k,V_k)_{k\ge1})$. Let $T_1$ be
the stopping time built from $B$ and $(U_1,V_1)$. By Theorem~\ref{thm:strongmarkov}, $B^{(1)}_t=B_{T_1+t}-B_{T_1}$ is a Brownian motion independent of
$\mathcal G_{T_1}$, which contains $\mathcal G_0$; build $T_2-T_1$ from $B^{(1)}$ and $(U_2,V_2)$ in the same way, and so on. By construction the pairs $(T_n-T_{n-1},
B_{T_n}-B_{T_{n-1}})$ are independent, each with the law of $(\tau,B_\tau)$ in Theorem~\ref{thm:skorokhod}.
\end{proof}

\section{Donsker's theorem}

Let $X_1,X_2,\dots$ be i.i.d.\ with mean $0$ and variance $1$, and $S_n=X_1+\dots+X_n$. Let $S(t)$ be the linear interpolation of the walk: $S(k)=S_k$ at integers, linear
in between. For each $n$, the random function
\[
X^n(t)=\frac{S(nt)}{\sqrt n},\qquad0\le t\le1,
\]
is a random element of $C$.

\begin{theorem}[Donsker's invariance principle, 1951]\label{thm:donsker}
$X^n\Rightarrow B$ in $C[0,1]$, where $B$ is a Brownian motion on $[0,1]$.
\end{theorem}
\begin{proof}
By Corollary~\ref{cor:embedwalk}, we may assume that $S_k=B_{T_k}$ for a Brownian motion $B$ and stopping times $T_k$ with i.i.d.\ increments of mean $1$. Let
$W^n(t)=B_{nt}/\sqrt n$; by scaling it is a Brownian motion, so it suffices, by Lemma~\ref{lem:together}, to show that $\|X^n-W^n\|\to0$ in probability.

By the strong law, $T_k/k\to1$ almost surely, and hence $\delta_n=\frac1n\max_{k\le n+1}|T_k-k|\to0$ almost surely. (Given $\eta>0$, $|T_k-k|\le\eta k$ for $k\ge k_0(\omega)$, so
$\delta_n\le\frac1n\max_{k<k_0}|T_k-k|+\eta\frac{n+1}n$.)

For $t\in[\frac kn,\frac{k+1}n]$, $X^n(t)$ lies between $S_k/\sqrt n$ and $S_{k+1}/\sqrt n$, so
\[
|X^n(t)-W^n(t)|\le\frac1{\sqrt n}\max\bigl(|B_{T_k}-B_{nt}|,|B_{T_{k+1}}-B_{nt}|\bigr).
\]
Here $|T_k-nt|\le|T_k-k|+1$ and the same for $k+1$, so on the event $\{\delta_n\le\delta\}$, with $n\ge1/\delta$, both times differ from $nt$ by at most $2\delta n$, and all
times involved are at most $2n$. Hence on that event
\[
\|X^n-W^n\|\le\sup\Bigl\{\frac{|B_r-B_s|}{\sqrt n}:r,s\le2n,\ |r-s|\le2\delta n\Bigr\},
\]
and the right side has the law of $w(2\delta)=\sup\{|B_r-B_s|:r,s\le2,|r-s|\le2\delta\}$, by scaling. So
\[
P(\|X^n-W^n\|>\varepsilon)\le P(\delta_n>\delta)+P(w(2\delta)>\varepsilon).
\]
The first term tends to $0$ as $n\to\infty$. The second tends to $0$ as $\delta\to0$, because Brownian paths are uniformly continuous on $[0,2]$, so $w(2\delta)\to0$ surely.
\end{proof}

\begin{figure}[ht]
\centering\fig{ch03-donsker}
\caption{One random walk with steps uniform on $[-\sqrt3,\sqrt3]$ (variance $1$), drawn as $X^n$ for $n=16$, $256$ and $65\,536$: the same walk, seen from further and
further away. The steps disappear and a Brownian path emerges (Theorem~\ref{thm:donsker}).}
\label{fig:donsker}
\end{figure}

Every functional $h:C\to\R$ continuous at almost every Brownian path now yields a limit theorem, valid for all walks at once, with the limit law of $h(B)$. And the
computation can be done for whichever walk is easiest, or for Brownian motion itself.

\begin{corollary}[the maximum]\label{cor:max}
For every walk with mean $0$ and variance $\sigma^2>0$, and $x\ge0$, $P(\max_{k\le n}S_k\ge x\sigma\sqrt n)\to P(|Z|\ge x)$ with $Z\sim N(0,1)$.
\end{corollary}
\begin{proof}
Apply Theorem~\ref{thm:donsker} to $X_k/\sigma$. The map $h(f)=\max_{t\le1}f(t)$ is $1$-Lipschitz on $C$, and
\[
h(X^n)=\max_{k\le n}S_k/(\sigma\sqrt n),
\]
since the maximum of a piecewise linear function is attained at a vertex. So $h(X^n)\Rightarrow\max_{t\le1}B_t$, which has the law of $|B_1|$ (Tome~3, Theorem~11.12). That law has a
density, so every point is a continuity point of its distribution function.
\end{proof}

The arcsine law of the time spent positive is subtler, because the functional is not continuous. Let
\[
h(f)=\lambda\{t\in[0,1]:f(t)>0\},
\]
with $\lambda$ the Lebesgue measure.

\begin{lemma}\label{lem:occupcont}
$h$ is measurable and continuous at every $f$ with $\lambda\{t:f(t)=0\}=0$. Brownian motion satisfies $\lambda\{t\le1:B_t=0\}=0$ almost surely.
\end{lemma}
\begin{proof}
$h(f)=\int_0^1\mathbf 1_{f(t)>0}\dd t$, and $(f,t)\mapsto\mathbf 1_{f(t)>0}$ is Borel on $C\times[0,1]$ (the evaluation $(f,t)\mapsto f(t)$ is continuous), so $h$ is measurable by Fubini. If
$f_n\to f$ uniformly, then $\mathbf 1_{f_n(t)>0}\to\mathbf 1_{f(t)>0}$ at every $t$ with $f(t)\ne0$, so by dominated convergence $h(f_n)\to h(f)$ when $\lambda\{f=0\}=0$. For Brownian
motion, Fubini gives $\mathbb E\lambda\{t:B_t=0\}=\int_0^1P(B_t=0)\dd t=0$.
\end{proof}

For the simple random walk ($X_k=\pm1$), $h(X^{2n})$ is the fraction of the $2n$ edges of the path lying above the axis (an edge from $S_{k-1}$ to $S_k$ lies above
when $S_{k-1}+S_k>0$). Its law is a gem of combinatorics.

\begin{lemma}[Chung--Feller]\label{lem:chungfeller}
For the simple random walk, the probability $b_{2k,2n}$ that exactly $2k$ of the first $2n$ edges lie above the axis is $u_{2k}u_{2n-2k}$, for $0\le k\le n$.
\end{lemma}
\begin{proof}
The number of edges above is even, because the path is above between consecutive returns to $0$, and these occur at even times. We argue by induction on $n$. For
$k=n$, all edges are above exactly when $S_1,\dots,S_{2n}\ge0$, which has probability $u_{2n}=u_{2n}u_0$ by Tome~3, Theorem~7.3; $k=0$ is the same by symmetry. This
covers $n=1$. Let $1\le k\le n-1$. Then the path has edges on both sides, so it returns to $0$; let $2r$ be the first return time, $P(T=2r)=f_{2r}$, and the first excursion
is above or below with probability $\frac12$ each, independently of $T$. After time $2r$ the walk starts afresh. If the first excursion is above, it contributes $2r$ edges
above, so $2r\le2k$; if below, it contributes none, so $2n-2r\ge2k$. Hence
\[
b_{2k,2n}=\frac12\sum_{r=1}^kf_{2r}b_{2k-2r,2n-2r}+\frac12\sum_{r=1}^{n-k}f_{2r}b_{2k,2n-2r}.
\]
By induction, $b_{2k-2r,2n-2r}=u_{2k-2r}u_{2n-2k}$ and $b_{2k,2n-2r}=u_{2k}u_{2n-2r-2k}$. The first-return decomposition (Tome~3, Lemma~7.2)
\[
\textstyle\sum_{r\le m}f_{2r}u_{2m-2r}=u_{2m},
\]
used with $m=k$ and $m=n-k$, turns the two sums into $\frac12u_{2n-2k}u_{2k}+\frac12u_{2k}u_{2n-2k}$.
\end{proof}

\begin{theorem}[L\'evy's arcsine law, 1939; Erd\H{o}s--Kac, 1947]\label{thm:arcsine}
For every walk with mean $0$ and finite positive variance, the fraction of time $h(X^n)$ that the rescaled path spends above $0$ converges in distribution to the
arcsine law of density $\frac1{\pi\sqrt{x(1-x)}}$ on $(0,1)$; and $\lambda\{t\le1:B_t>0\}$ has the arcsine law.
\end{theorem}
\begin{proof}
By Theorem~\ref{thm:donsker}, Lemma~\ref{lem:occupcont} and Proposition~\ref{prop:contmap}, $h(X^n)\Rightarrow h(B)$ for every such walk. For the simple random walk along
even $n$, Lemma~\ref{lem:chungfeller} says that $h(X^{2n})=k/n$ with probability $u_{2k}u_{2n-2k}$, and Tome~3, Theorem~7.7 says exactly that
$P(a\le h(X^{2n})\le b)\to\frac2\pi(\arcsin\sqrt b-\arcsin\sqrt a)$ for $0<a<b<1$. The limit law of a weakly convergent sequence is unique and determined by these values, so
$h(B)$ has the arcsine law, and then so does the limit for every walk.
\end{proof}

\begin{think}
The arcsine density is largest near $0$ and $1$. So in a fair game of $1000$ rounds, what is more likely: that the leader changes often and each player leads about half the
time, or that one player leads for more than $90\%$ of the time? Compute $P(h(B)>0.9)+P(h(B)<0.1)$, and then guess how often the lead changes in $1000$ rounds of
coin tossing. The notebook checks.
\end{think}

\begin{remark}[Kolmogorov--Smirnov]
The same philosophy governs goodness-of-fit tests. If $U_1,\dots,U_n$ are i.i.d.\ uniform on $[0,1]$ and $F_n$ is their empirical distribution function, then
$\sqrt n(F_n(t)-t)$ converges in distribution, as a random function, to the \emph{Brownian bridge} $B_t-tB_1$ (Donsker, 1952), and so
$\sqrt n\sup_t|F_n(t)-t|$ converges to $\sup_t|B_t-tB_1|$, whose distribution function Kolmogorov had found in 1933:
$1-2\sum_{k\ge1}(-1)^{k-1}e^{-2k^2x^2}$. The proof needs a space of functions with jumps, which we do not develop; see Billingsley's book in the Sources. The notebook tests
random number generators with this statistic.
\end{remark}

\section{The law of the iterated logarithm}

Let $h(t)=\sqrt{2t\log\log t}$ for $t\ge3$. We use the Gaussian tail bounds, for $x>0$ and $Z\sim N(0,1)$,
\begin{equation}\label{eq:gausstail}
\frac x{1+x^2}\frac{e^{-x^2/2}}{\sqrt{2\pi}}\le P(Z\ge x)\le\frac12e^{-x^2/2},
\end{equation}
the left one from $\frac{\dd}{\dd x}\bigl(-\frac x{1+x^2}e^{-x^2/2}\bigr)\le e^{-x^2/2}$ (Exercise~\ref{ex:gausstail}), the right one from $P(Z\ge x)\le\mathbb E e^{x(Z-x)}\mathbf 1_{Z\ge x}$
and the shift $Z\mapsto Z+x$ of the Gaussian density, which turns that expectation into $e^{-x^2/2}P(Z\ge0)$.

\begin{theorem}[law of the iterated logarithm for Brownian motion]\label{thm:lilbm}
$\displaystyle\limsup_{t\to\infty}\frac{B_t}{h(t)}=1$ almost surely, and the same holds along the integers.
\end{theorem}
\begin{proof}
\emph{Upper bound.} Let $\varepsilon>0$, $1<q<(1+\varepsilon)^2$, and $t_k=q^k$. By the reflection principle and~\eqref{eq:gausstail},
\[
P\Bigl(\max_{s\le t_{k+1}}B_s\ge(1+\varepsilon)h(t_k)\Bigr)=2P\bigl(B_{t_{k+1}}\ge(1+\varepsilon)h(t_k)\bigr)\le\exp\Bigl(-\frac{(1+\varepsilon)^2h(t_k)^2}{2t_{k+1}}\Bigr)
=\bigl(k\log q\bigr)^{-(1+\varepsilon)^2/q}.
\]
The exponent exceeds $1$, so the series converges and, by Borel--Cantelli, for all large $k$ and all $t\in[t_k,t_{k+1}]$, $B_t<(1+\varepsilon)h(t_k)\le(1+\varepsilon)h(t)$. So
$\limsup B_t/h(t)\le1+\varepsilon$.

\emph{Lower bound.} Let $\varepsilon\in(0,1)$, $q>1$ large, and $t_k=\lceil q^k\rceil$, integers. The increments $D_k=B_{t_k}-B_{t_{k-1}}$ are independent, $N(0,t_k-t_{k-1})$. Let
$x_k=(1-\varepsilon)h(t_k)/\sqrt{t_k-t_{k-1}}$. Since $t_k-t_{k-1}\ge t_k(1-\frac2q)$ for large $k$,
\[
\frac{x_k^2}2\le\frac{(1-\varepsilon)^2}{1-2/q}\log\log t_k=\gamma\log\log t_k,\qquad\gamma=\frac{(1-\varepsilon)^2}{1-2/q},
\]
and we choose $q$ so large that $\gamma<1$. By~\eqref{eq:gausstail}, $P(D_k\ge(1-\varepsilon)h(t_k))\ge\frac c{x_k}(\log t_k)^{-\gamma}\ge\frac{c'}{\sqrt{\log k}\,k^\gamma}$, whose sum diverges. By the
second Borel--Cantelli lemma, $D_k\ge(1-\varepsilon)h(t_k)$ infinitely often. By the upper bound applied to $-B$, $B_{t_{k-1}}\ge-2h(t_{k-1})$ for large $k$, and $h(t_{k-1})\le
\frac2{\sqrt q}h(t_k)$ for large $k$. So infinitely often
\[
B_{t_k}=D_k+B_{t_{k-1}}\ge\Bigl(1-\varepsilon-\frac4{\sqrt q}\Bigr)h(t_k).
\]
Since $\varepsilon$ is arbitrary and $q$ can be taken arbitrarily large,
\[
\limsup_kB_{t_k}/h(t_k)\ge1,
\]
along integer times.
\end{proof}

\begin{theorem}[Hartman--Wintner, 1941]\label{thm:hartmanwintner}
Let $X_k$ be i.i.d.\ with mean $0$ and variance $\sigma^2>0$. Then
\[
\limsup_{n\to\infty}\frac{S_n}{\sigma\sqrt{2n\log\log n}}=1\quad\text{and}\quad\liminf_{n\to\infty}\frac{S_n}{\sigma\sqrt{2n\log\log n}}=-1\quad\text{almost surely.}
\]
\end{theorem}
\begin{proof}
We may take $\sigma=1$ and, by Corollary~\ref{cor:embedwalk}, $S_n=B_{T_n}$. The $\liminf$ follows from the $\limsup$ applied to $-X_k$. By Theorem~\ref{thm:lilbm} along integers, it
suffices to show $|B_{T_n}-B_n|/h(n)\to0$ almost surely.

Let $0<\varepsilon\le\frac14$ and $r=1+\varepsilon$. By the strong law, $|T_n-n|\le\varepsilon n$ for large $n$. For $n\in[r^k,r^{k+1}]$, both $n$ and $T_n$ lie in $[r^{k-2},r^{k+2}]$,
because $(1-\varepsilon)r^2\ge1$ and $(1+\varepsilon)\le r$. So $|B_{T_n}-B_n|\le2M_k$, where $M_k=\max_{0\le u\le L_k}|B_{r^{k-2}+u}-B_{r^{k-2}}|$ and $L_k=r^{k+2}-r^{k-2}\le
r^k\bigl(r^2-r^{-2}\bigr)\le5\varepsilon r^k$. By the Markov property and the reflection principle applied to $B$ and $-B$, $P(M_k\ge x)\le4P(B_{L_k}\ge x)\le2e^{-x^2/2L_k}$. With
$x=\delta h(r^k)$ and $\delta^2=20\varepsilon$,
\[
P\bigl(M_k\ge\delta h(r^k)\bigr)\le2\exp\Bigl(-\frac{20\varepsilon\cdot2r^k\log\log r^k}{2\cdot5\varepsilon r^k}\Bigr)=2(k\log r)^{-4},
\]
summable. So for large $k$, $M_k<\delta h(r^k)\le\delta h(n)$, and $\limsup|B_{T_n}-B_n|/h(n)\le2\sqrt{20\varepsilon}$. Let $\varepsilon\to0$.
\end{proof}

\begin{figure}[ht]
\centering\fig{ch03-lil}
\caption{Four coin-tossing walks, as $S_n/\sqrt n$ against $n$ on a logarithmic scale, with the envelopes $\pm\sqrt{2\log\log n}$ of Theorem~\ref{thm:hartmanwintner}. The
walks cross the envelope from time to time, but at the scale of the theorem, $\log\log n$, a million steps is still very short: $\sqrt{2\log\log10^6}=2.29$.}
\label{fig:lil}
\end{figure}

\begin{beauty}
Donsker's theorem is the central limit theorem for whole paths, and the reason it holds is almost physical: Brownian motion contains every random walk with finite variance,
stopped at the right times, and on a large scale the stopping times tick like a clock, by the strong law. So one universal object answers infinitely many questions, and each
answer can be computed wherever it is easiest: by counting lattice paths, by reflecting Brownian motion, or by a computer.
\end{beauty}

\section{Exercises}

\begin{exercise}[$\star$]\label{ex:gausstail}
Prove the Gaussian tail bounds~\eqref{eq:gausstail}. (For the lower bound, differentiate $-\frac x{1+x^2}e^{-x^2/2}$.)
\end{exercise}

\begin{exercise}[$\star$]\label{ex:multinomial}
Let $(N_1,\dots,N_d)$ be multinomial with $n$ trials and probabilities $p_1,\dots,p_d$. Show with the Cram\'er--Wold device that $\bigl((N_i-np_i)/\sqrt n\bigr)_i\Rightarrow N(0,\Sigma)$
with $\Sigma=\operatorname{diag}(p)-pp^{\mathsf T}$. Deduce that Pearson's statistic $\sum_i(N_i-np_i)^2/(np_i)$ converges in distribution to $\chi^2_{d-1}$.
\end{exercise}

\begin{exercise}[$\star$]\label{ex:bridge}
Show that $B_t-tB_1$, $0\le t\le1$, is a centred Gaussian process with covariance $s(1-t)$ for $s\le t$, independent of $B_1$. Deduce from Donsker's theorem that
$X^n(t)-tX^n(1)\Rightarrow B_t-tB_1$ in $C[0,1]$.
\end{exercise}

\begin{exercise}[$\star\star$]\label{ex:exitvar}
For the exit time $\tau$ of $(-1,1)$ by Brownian motion, show that $B_t^4-6tB_t^2+3t^2$ is a martingale and deduce $\mathbb E\tau^2=\frac53$, so $\operatorname{Var}\tau=\frac23$. (These
are the times between steps when the simple random walk is embedded in Brownian motion.)
\end{exercise}

\begin{exercise}[$\star\star$]\label{ex:clusterset}
Under the hypotheses of Theorem~\ref{thm:hartmanwintner}, show that almost surely the set of limit points of $S_n/(\sigma\sqrt{2n\log\log n})$ is the whole interval $[-1,1]$.
(Show first that $X_n/\sqrt n\to0$ almost surely.)
\end{exercise}

\begin{exercise}[$\star\star$]\label{ex:lastzero}
Let $g=\sup\{t\le1:B_t=0\}$. Show that $g$ has the arcsine law by computing $P(g<s)=P(B_u\ne0\text{ for }u\in[s,1])$ with the Markov property at time $s$ and the reflection
principle: $P(g<s)=\mathbb E\bigl[P(|Z|\sqrt{1-s}<|B_s|)\bigr]$ with $Z$ independent of $B_s$, then evaluate the probability that a planar Gaussian vector lies in a cone.
\end{exercise}

\begin{exercise}[$\star\star\star$]\label{ex:modulus}
Show that Brownian paths are not H\"older continuous of order $\frac12$: almost surely, for every $c<\sqrt2$, there are infinitely many $m$ such that
$\max_{k\le2^m}|B_{k2^{-m}}-B_{(k-1)2^{-m}}|>c\sqrt{m2^{-m}\log2}$. (Use independence of the increments, the lower bound in~\eqref{eq:gausstail}, $1-y\le e^{-y}$, and
Borel--Cantelli.)
\end{exercise}

\section{Hints}

\begin{hint}{ex:gausstail}
The derivative of $-\frac x{1+x^2}e^{-x^2/2}$ is $\bigl(1-\frac2{(1+x^2)^2}\bigr)e^{-x^2/2}\le e^{-x^2/2}$; integrate from $x$ to $\infty$. For the upper bound, $\mathbf 1_{Z\ge x}\le e^{x(Z-x)}\mathbf 1_{Z\ge x}$
and compute $\mathbb E[e^{xZ}\mathbf 1_{Z\ge x}]$ by completing the square.
\end{hint}
\begin{hint}{ex:multinomial}
$N=\sum_{k\le n}\xi_k$ with $\xi_k$ i.i.d.\ unit vectors $e_i$ with probability $p_i$. For the $\chi^2$ part, write the statistic as $|D^{-1/2}Y|^2$ and diagonalise
$D^{-1/2}\Sigma D^{-1/2}$, a projection of rank $d-1$.
\end{hint}
\begin{hint}{ex:bridge}
Compute covariances; a Gaussian vector with uncorrelated blocks has independent blocks. The map $f\mapsto f-tf(1)$ is continuous on $C$.
\end{hint}
\begin{hint}{ex:exitvar}
Expand $\mathbb E[(B_s+N)^4\mid\mathcal F_s]$. Apply optional stopping at $\tau\wedge K$ (justify as in Lemma~\ref{lem:ostbm}) and let $K\to\infty$; $\mathbb E\tau=1$.
\end{hint}
\begin{hint}{ex:clusterset}
$\sum P(X_n^2>\varepsilon^2n)\le\mathbb E X^2/\varepsilon^2$. Then consecutive values of $S_n/h(n)$ differ by a quantity tending to $0$, and the sequence goes from near $1$ to near $-1$
infinitely often.
\end{hint}
\begin{hint}{ex:lastzero}
Given $B_s=x$, the probability of no zero in $[s,1]$ is $P(\max_{u\le1-s}W_u<|x|)=P(|W_{1-s}|<|x|)$. The pair $(B_s/\sqrt s,Z)$ is a standard planar Gaussian, rotation invariant.
\end{hint}
\begin{hint}{ex:modulus}
$P(\text{all }2^m\text{ increments}\le c\sqrt{m2^{-m}\log2})=\bigl(1-P(|Z|>c\sqrt{m\log2})\bigr)^{2^m}\le\exp\bigl(-2^mP(|Z|>c\sqrt{m\log2})\bigr)$.
\end{hint}

\section{Solutions}

\begin{solution}{ex:gausstail}
\idea Both bounds compare the Gaussian tail integral with something explicit: an exact antiderivative of a slightly smaller function for the lower bound, an exponential
moment for the upper bound.
\writeup Let $\phi(x)=e^{-x^2/2}/\sqrt{2\pi}$ and $G(x)=-\frac x{1+x^2}\phi(x)$. Then $G'(x)=\bigl(\frac{x^2}{1+x^2}-\frac{1-x^2}{(1+x^2)^2}\bigr)\phi(x)=\bigl(1-\frac2{(1+x^2)^2}\bigr)\phi(x)\le\phi(x)$. Integrating
from $x$ to $\infty$, where $G\to0$: $\frac x{1+x^2}\phi(x)=-G(x)=\int_x^\infty G'\le\int_x^\infty\phi=P(Z\ge x)$. For the upper bound, $P(Z\ge x)\le\mathbb E[e^{x(Z-x)}\mathbf 1_{Z\ge x}]=
e^{-x^2}\int_x^\infty e^{xz}\phi(z)\dd z=e^{-x^2/2}\int_x^\infty\phi(z-x)\dd z=\frac12e^{-x^2/2}$.
\end{solution}

\begin{solution}{ex:multinomial}
\idea A multinomial vector is a sum of $n$ i.i.d.\ indicator vectors, so this is a multivariate central limit theorem, which Cram\'er--Wold reduces to the real one. The
$\chi^2$ limit is then a continuous mapping.
\writeup Write $N=\sum_{k\le n}\xi_k$ with $\xi_k$ i.i.d., $P(\xi_k=e_i)=p_i$. Then $\mathbb E\xi_k=p$ and $\operatorname{Cov}\xi_k=\mathbb E[\xi\xi^{\mathsf T}]-pp^{\mathsf T}=\operatorname{diag}(p)-pp^{\mathsf T}=\Sigma$. For each
$t$, $\langle t,\xi_k\rangle$ are i.i.d.\ real with variance $t^{\mathsf T}\Sigma t$, so by Tome~3's CLT $\langle t,(N-np)/\sqrt n\rangle\Rightarrow N(0,t^{\mathsf T}\Sigma t)$, which is the law of
$\langle t,Y\rangle$ for $Y\sim N(0,\Sigma)$. By Corollary~\ref{cor:cramerwold}, $(N-np)/\sqrt n\Rightarrow Y$. Pearson's statistic is $g\bigl((N-np)/\sqrt n\bigr)$ with $g(y)=|D^{-1/2}y|^2$,
$D=\operatorname{diag}(p)$, continuous, so it converges to $|D^{-1/2}Y|^2$. Now $D^{-1/2}Y\sim N(0,P)$ with $P=I-\sqrt p\sqrt p^{\mathsf T}$, where $\sqrt p$ is the unit vector $(\sqrt{p_i})$:
$P$ is the orthogonal projection onto $\sqrt p^\perp$, of rank $d-1$. In an orthonormal basis adapted to it, $D^{-1/2}Y$ has $d-1$ independent standard coordinates and one zero coordinate,
so $|D^{-1/2}Y|^2\sim\chi^2_{d-1}$.
\end{solution}

\begin{solution}{ex:bridge}
\idea Everything is Gaussian, so independence is a covariance computation, and the convergence is a continuous mapping of Donsker's theorem.
\writeup $(B_{t_1}-t_1B_1,\dots,B_{t_k}-t_kB_1,B_1)$ is a linear image of a Gaussian vector, hence Gaussian (Tome~3, Theorem~6.15). For $s\le t$,
$\operatorname{Cov}(B_s-sB_1,B_t-tB_1)=s-st-st+st=s(1-t)$, and $\operatorname{Cov}(B_t-tB_1,B_1)=t-t=0$. Uncorrelated blocks of a Gaussian vector are independent, so the bridge is independent
of $B_1$ (Lemma~\ref{lem:pisuffice} extends this from finitely many times to the whole process). The map $\Phi(f)(t)=f(t)-tf(1)$ is linear and satisfies $\|\Phi f\|\le2\|f\|$, so it is
continuous on $C$, and $\Phi(X^n)\Rightarrow\Phi(B)$ by Proposition~\ref{prop:contmap}.
\end{solution}

\begin{solution}{ex:exitvar}
\idea Fourth moments of $\tau$-stopped Brownian motion come from a polynomial martingale of degree four, the next one in the family $B_t$, $B_t^2-t$; its coefficients are forced by
the Gaussian moments.
\writeup With $N=B_t-B_s\sim N(0,u)$, $u=t-s$, independent of $\mathcal F_s$, $\mathbb E[B_t^4\mid\mathcal F_s]=B_s^4+6B_s^2u+3u^2$ and $\mathbb E[B_t^2\mid\mathcal F_s]=B_s^2+u$. So
\[
\begin{gathered}
\mathbb E[B_t^4-6tB_t^2+3t^2\mid\mathcal F_s]=B_s^4+6B_s^2u+3u^2-6t(B_s^2+u)+3t^2\qq=B_s^4-6sB_s^2+3s^2,
\end{gathered}
\]
since $6u-6t=-6s$ and $3u^2-6tu+3t^2=3s^2$. Optional stopping at $\tau\wedge K$ (the dyadic approximation of Lemma~\ref{lem:ostbm} works, with the domination
$\max_{s\le K+1}B_s^4+(K+1)^2\max_sB_s^2$, integrable) gives $\mathbb E B_{\tau\wedge K}^4+3\mathbb E(\tau\wedge K)^2=6\mathbb E[(\tau\wedge K)B_{\tau\wedge K}^2]$. Since $|B_{\tau\wedge K}|\le1$,
$3\mathbb E(\tau\wedge K)^2\le6\mathbb E\tau=6$, so $\mathbb E\tau^2<\infty$, and dominated convergence as $K\to\infty$ gives $1+3\mathbb E\tau^2=6\mathbb E\tau=6$, since $B_\tau^2=1$. So $\mathbb E\tau^2=\frac53$
and $\operatorname{Var}\tau=\frac53-1=\frac23$.
\end{solution}

\begin{solution}{ex:clusterset}
\idea A sequence whose consecutive gaps tend to $0$ and which goes from near $+1$ to near $-1$ infinitely often must pass near every point in between.
\writeup Let $h_n=\sigma\sqrt{2n\log\log n}$ and $y_n=S_n/h_n$. Since $\sum_nP(X_n^2>\varepsilon^2n)=\sum_nP(X^2/\varepsilon^2>n)\le\mathbb E X^2/\varepsilon^2<\infty$, Borel--Cantelli gives
$|X_n|\le\varepsilon\sqrt n$ eventually, so $X_n/\sqrt n\to0$ and $X_{n+1}/h_{n+1}\to0$. Moreover $|S_n|/h_n$ is bounded (it tends to at most $1$ in $\limsup$), and $h_n/h_{n+1}\to1$, so
$y_{n+1}-y_n=\frac{X_{n+1}}{h_{n+1}}+y_n\bigl(\frac{h_n}{h_{n+1}}-1\bigr)\to0$. By Theorem~\ref{thm:hartmanwintner}, $y_n>1-\eta$ and $y_n<-1+\eta$ each happen infinitely often. Given
$c\in(-1,1)$ and $\eta>0$, between a time when $y_n<c$ and a later time when $y_n>c$ the sequence crosses $c$ with steps eventually smaller than $\eta$, so it lands in $(c-\eta,c+\eta)$.
Hence every $c\in[-1,1]$ is a limit point (the endpoints by the theorem itself), and no point outside $[-1,1]$ is one.
\end{solution}

\begin{solution}{ex:lastzero}
\idea ``No zero after time $s$'' is a statement about the Brownian motion started afresh at time $s$ from the point $B_s$, which the reflection principle answers; what is left is
the probability that one Gaussian is larger than another, which is an angle.
\writeup By the Markov property at $s$, $W_u=B_{s+u}-B_s$ is a Brownian motion independent of $B_s$. On $\{B_s=x\}$ with $x>0$, there is no zero in $[s,1]$ iff $\min_{u\le1-s}W_u>-x$,
of probability $P(|W_{1-s}|<x)$ by the reflection principle and symmetry; likewise for $x<0$. So $P(g<s)=P(\sqrt{1-s}\,|Z|<\sqrt s\,|Z'|)$ with $Z,Z'$ independent standard
Gaussians. The vector $(Z',Z)$ has a rotation-invariant law; the event $|Z|<\sqrt{\frac s{1-s}}|Z'|$ is a double cone of half-angle $\theta=\arctan\sqrt{\frac s{1-s}}=\arcsin\sqrt s$ around the
$Z'$-axis, so its probability is $\frac{4\theta}{2\pi}=\frac2\pi\arcsin\sqrt s$. This is the arcsine distribution function.
\end{solution}

\begin{solution}{ex:modulus}
\idea On level $m$ there are $2^m$ independent increments of variance $2^{-m}$; the largest of $N$ independent standard Gaussians is about $\sqrt{2\log N}$. So the largest increment is
about $\sqrt{2\log(2^m)\,2^{-m}}$, and anything smaller by a constant factor fails with a probability small enough to sum.
\writeup Let $x_m=c\sqrt{m\log2}$. The normalised increments $2^{m/2}(B_{k2^{-m}}-B_{(k-1)2^{-m}})$, $k\le2^m$, are i.i.d.\ standard Gaussians. So the probability that all of them are at most
$x_m$ in absolute value is
\[
\bigl(1-P(|Z|>x_m)\bigr)^{2^m}\le\exp\bigl(-2^mP(|Z|>x_m)\bigr).
\]
By~\eqref{eq:gausstail}, $P(|Z|>x_m)\ge\frac{c_1}{\sqrt m}e^{-x_m^2/2}=\frac{c_1}{\sqrt m}2^{-c^2m/2}$, so $2^mP(|Z|>x_m)\ge\frac{c_1}{\sqrt m}2^{m(1-c^2/2)}$, which grows exponentially in
$m$ since $c^2<2$. The probabilities $\exp(-\cdots)$ are therefore summable, and by Borel--Cantelli, for all large $m$ some increment exceeds $x_m2^{-m/2}=c\sqrt{m2^{-m}\log2}$. Taking
$c=\sqrt2-\frac1j$ for $j=1,2,\dots$ handles all $c<\sqrt2$ at once. In particular $|B_t-B_s|\le C|t-s|^{1/2}$ fails for every $C$: with $|t-s|=2^{-m}$ the ratio is at least
$c\sqrt{m\log2}\to\infty$.
\end{solution}

\section*{Sources and further reading}

\begin{itemize}
\item \emph{History.} P.~Erd\H{o}s and M.~Kac, ``On certain limit theorems of the theory of probability'', \emph{Bull.\ Amer.\ Math.\ Soc.}\ 52 (1946), 292--302, and ``On the number of
positive sums of independent random variables'', \emph{Bull.\ Amer.\ Math.\ Soc.}\ 53 (1947), 1011--1020. M.~D.~Donsker, ``An invariance principle for certain probability limit theorems'',
\emph{Mem.\ Amer.\ Math.\ Soc.}\ 6 (1951); ``Justification and extension of Doob's heuristic approach to the Kolmogorov--Smirnov theorems'', \emph{Ann.\ Math.\ Statist.}\ 23 (1952),
277--281. J.~L.~Doob, ``Heuristic approach to the Kolmogorov--Smirnov theorems'', \emph{Ann.\ Math.\ Statist.}\ 20 (1949), 393--403. A.~N.~Kolmogorov, ``Sulla determinazione empirica di una
legge di distribuzione'', \emph{Giorn.\ Ist.\ Ital.\ Attuari} 4 (1933), 83--91. A.~V.~Skorokhod, \emph{Studies in the Theory of Random Processes}, Kiev, 1961 (English translation,
Addison-Wesley, 1965). A.~Khintchine, ``\"Uber einen Satz der Wahrscheinlichkeitsrechnung'', \emph{Fund.\ Math.}\ 6 (1924), 9--20. P.~Hartman and A.~Wintner, ``On the law of the iterated
logarithm'', \emph{Amer.\ J.\ Math.}\ 63 (1941), 169--176. P.~L\'evy, ``Sur certains processus stochastiques homog\`enes'', \emph{Compositio Math.}\ 7 (1939), 283--339 (the arcsine laws).
K.~L.~Chung and W.~Feller, ``On fluctuations in coin-tossing'', \emph{Proc.\ Nat.\ Acad.\ Sci.\ USA} 35 (1949), 605--608.
\item R.~Durrett, \emph{Probability: Theory and Examples}, 5th ed., Cambridge University Press, 2019, Section~3.3 (L\'evy's theorem) and Chapters~7--8 (Brownian motion, Skorokhod embedding, Donsker, the
law of the iterated logarithm). The route followed here.
\item P.~Billingsley, \emph{Convergence of Probability Measures}, 2nd ed., Wiley, 1999. Tightness in $C$ and in the space $D$ of functions with jumps; empirical processes and
Kolmogorov--Smirnov.
\item W.~Feller, \emph{An Introduction to Probability Theory and Its Applications}, vol.~I, 3rd ed., Wiley, 1968, Chapter~III. The Chung--Feller theorem and the surprising
fluctuations of coin tossing.
\item P.~M\"orters and Y.~Peres, \emph{Brownian Motion}, Cambridge University Press, 2010, Chapters~2 and~5. The strong Markov property, embeddings and the law of the iterated logarithm.
\end{itemize}
